% =========================================================
% Proof Stub: Characterization of Limsup via Subsequences
% Source: volume-iii/analysis/sequences/notes/asymptotics/notes-limsup-liminf.tex
% Mode: standard
% =========================================================

% *** HARDER PROOF — attempt after foundational results settled ***

\clearpage
\phantomsection
\hypertarget{proof-RA-ABB-T-limsup-subseq-char}{}

\subsubsection[Characterization of Limsup via Subsequences]{Proof --- Characterization of Limsup via Subsequences}

\bigskip

\noindent
\textbf{Source.}
Abbott, \emph{Understanding Analysis}; see \texttt{notes-limsup-liminf.tex}.

\vspace{0.75em}

\noindent
\textbf{Goal.}
$L = \limsup x_n$ iff (i) there is a subsequence $(x_{n_k}) \to L$ and (ii) for all $\varepsilon > 0$, eventually $x_n < L + \varepsilon$.

\vspace{0.75em}

\noindent
\textbf{Logical form.}
\[
  L = \limsup x_n \Longleftrightarrow (\exists (x_{n_k}) \to L) \wedge (\forall \varepsilon > 0,\; \exists N,\; \forall n \ge N,\; x_n < L + \varepsilon).
\]

\vspace{0.75em}

\noindent
\textbf{Proof strategy.}
Use the tail-supremum characterization: $s_n \to L$ means $s_n$ gets close to $L$, giving a convergent subsequence near $L$ from above and an eventual upper bound near $L$.

\vspace{0.75em}

\noindent
\textbf{Proof.}
\begin{proof}
\Claim $L = \limsup x_n$ iff (i) there is a subsequence $(x_{n_k}) \to L$ and (ii) for all $\varepsilon > 0$, eventually $x_n < L + \varepsilon$.

\Given 

\Goal 

\Strategy Use the tail-supremum characterization: $s_n \to L$ means $s_n$ gets close to $L$, giving a convergent subsequence near $L$ from above and an eventual upper bound near $L$.

\medskip

% --- (⇒) s_n → L: for each k, s_{n_k} - 1/k < s_{n_k} = sup_{j≥n_k} x_j, so ∃ x_{n_k}' > s_{n_k} - 1/k → L ---
% --- this subsequence converges to L ---
% --- also s_n → L gives eventual upper bound L + ε ---
% --- (⇐) show s_n → L using both conditions ---

\end{proof}

\begin{remark*}[Proof shape]
Two-sided characterization: a subsequence touches $L$ from below; the upper bound condition prevents values far above $L$.
\end{remark*}

\begin{remark*}[Dependencies]
\begin{itemize}
  \item Tail suprema definition and monotonicity.
  \item Dense subsequence criterion.
  \item Squeeze theorem.
\end{itemize}
\end{remark*}
